% Options for packages loaded elsewhere
\PassOptionsToPackage{unicode}{hyperref}
\PassOptionsToPackage{hyphens}{url}
%
\documentclass[
  11pt,
]{article}
\usepackage{amsmath,amssymb}
\usepackage{iftex}
\ifPDFTeX
  \usepackage[T1]{fontenc}
  \usepackage[utf8]{inputenc}
  \usepackage{textcomp} % provide euro and other symbols
\else % if luatex or xetex
  \usepackage{unicode-math} % this also loads fontspec
  \defaultfontfeatures{Scale=MatchLowercase}
  \defaultfontfeatures[\rmfamily]{Ligatures=TeX,Scale=1}
\fi
\usepackage{lmodern}
\ifPDFTeX\else
  % xetex/luatex font selection
\fi
% Use upquote if available, for straight quotes in verbatim environments
\IfFileExists{upquote.sty}{\usepackage{upquote}}{}
\IfFileExists{microtype.sty}{% use microtype if available
  \usepackage[]{microtype}
  \UseMicrotypeSet[protrusion]{basicmath} % disable protrusion for tt fonts
}{}
\makeatletter
\@ifundefined{KOMAClassName}{% if non-KOMA class
  \IfFileExists{parskip.sty}{%
    \usepackage{parskip}
  }{% else
    \setlength{\parindent}{0pt}
    \setlength{\parskip}{6pt plus 2pt minus 1pt}}
}{% if KOMA class
  \KOMAoptions{parskip=half}}
\makeatother
\usepackage{xcolor}
\usepackage[margin=1in]{geometry}
\usepackage{longtable,booktabs,array}
\usepackage{calc} % for calculating minipage widths
% Correct order of tables after \paragraph or \subparagraph
\usepackage{etoolbox}
\makeatletter
\patchcmd\longtable{\par}{\if@noskipsec\mbox{}\fi\par}{}{}
\makeatother
% Allow footnotes in longtable head/foot
\IfFileExists{footnotehyper.sty}{\usepackage{footnotehyper}}{\usepackage{footnote}}
\makesavenoteenv{longtable}
\usepackage{graphicx}
\makeatletter
\def\maxwidth{\ifdim\Gin@nat@width>\linewidth\linewidth\else\Gin@nat@width\fi}
\def\maxheight{\ifdim\Gin@nat@height>\textheight\textheight\else\Gin@nat@height\fi}
\makeatother
% Scale images if necessary, so that they will not overflow the page
% margins by default, and it is still possible to overwrite the defaults
% using explicit options in \includegraphics[width, height, ...]{}
\setkeys{Gin}{width=\maxwidth,height=\maxheight,keepaspectratio}
% Set default figure placement to htbp
\makeatletter
\def\fps@figure{htbp}
\makeatother
\setlength{\emergencystretch}{3em} % prevent overfull lines
\providecommand{\tightlist}{%
  \setlength{\itemsep}{0pt}\setlength{\parskip}{0pt}}
\setcounter{secnumdepth}{-\maxdimen} % remove section numbering
\usepackage{amsmath,amssymb,mathtools}
\usepackage{booktabs}
\usepackage{array}
\renewcommand{\arraystretch}{1.25}
\allowdisplaybreaks
\setlength{\emergencystretch}{3em}
\usepackage{needspace}
\let\oldsubsection\subsection
\renewcommand{\subsection}{\needspace{6\baselineskip}\oldsubsection}
\ifLuaTeX
  \usepackage{selnolig}  % disable illegal ligatures
\fi
\IfFileExists{bookmark.sty}{\usepackage{bookmark}}{\usepackage{hyperref}}
\IfFileExists{xurl.sty}{\usepackage{xurl}}{} % add URL line breaks if available
\urlstyle{same}
\hypersetup{
  pdftitle={Stress and slip follow the lattice},
  hidelinks,
  pdfcreator={LaTeX via pandoc}}

\title{Stress and slip follow the lattice}
\usepackage{etoolbox}
\makeatletter
\providecommand{\subtitle}[1]{% add subtitle to \maketitle
  \apptocmd{\@title}{\par {\large #1 \par}}{}{}
}
\makeatother
\subtitle{Cosserat--KWC crystal plasticity at finite deformation
(production 7.5): statement, proof and numerical verification}
\author{}
\date{}

\begin{document}
\maketitle

\hypertarget{statement}{%
\section{1. Statement}\label{statement}}

\textbf{Claim.} With the elasticity tensor and the slip systems rotated
once by the initial orientation and never updated, the model's response
refers to a single lattice at every point and time. That lattice has the
current orientation \(R_t = RR_0\). Three things follow:

\begin{enumerate}
\def\labelenumi{\arabic{enumi}.}
\tightlist
\item
  The stress is the crystal's own elastic response, evaluated at the
  elastic stretch measured in the lattice frame and rotated by \(R_t\).
\item
  The resolved shear stresses act on, and the slip occurs along, the
  slip systems of the current lattice.
\item
  The lattice carrying (1) and (2) is the same lattice whose orientation
  enters the KWC curvature.
\end{enumerate}

All of this holds whatever made the lattice turn: rigid rotation, slip,
the grain-boundary spin, or any combination, in any order.

\hypertarget{setting}{%
\section{2. Setting}\label{setting}}

\textbf{Kinematics.}
\[F=R\,\bar U F_p,\qquad \bar U=R^TFF_p^{-1},\qquad R_t=RR_0,\qquad R(X,0)=I,\quad F_p(X,0)=I .\]
Here \(R\) is the relative Cosserat (lattice) rotation and \(R_0(X)\)
the initial orientation. Write \(F_e:=FF_p^{-1}=R_eU_e\) (polar
decomposition).

\textbf{Crystal data.} The data are rotated once, at setup, and never
updated:
\[\mathbb C=R_0\star\mathbb C^{c},\qquad s_\alpha=R_0s^c_\alpha,\qquad n_\alpha=R_0n^c_\alpha,\qquad (Q\star\mathbb C)_{ijkl}=Q_{ia}Q_{jb}Q_{kc}Q_{ld}\,\mathbb C_{abcd}.\]

\textbf{Plastic flow.} Slip plus the grain-boundary spin:
\[L_p=\dot F_pF_p^{-1}=\sum_\alpha\dot\gamma_\alpha\,s_\alpha\otimes n_\alpha+\dot w^*\,\widehat W .\]

\textbf{Energy.}
\[\psi=\psi_{el}(\bar U)+2\lambda\cdot a+2\mu_c|a|^2+\psi_K(K,\phi,\dots),\qquad a=\operatorname{axl}(\operatorname{skw}\bar U),\]
where the elastic part is written in the Green strain
\(\bar E=\tfrac12(\bar U^T\bar U-I)\) (plus a volumetric term in
\(J=\det\bar U\)). Its derivative is
\(\Pi=\partial\psi/\partial\bar U\).

\textbf{Stresses.}
\[P=R\,\Pi\,F_p^{-T},\qquad \tau=PF^T,\qquad M=\bar U^T\Pi,\qquad \tau_\alpha=M:(s_\alpha\otimes n_\alpha).\]

\textbf{Crystal-frame quantities.} For any reference-frame tensor \(A\),
write \(A^c:=R_0^TAR_0\). Let \(\psi^c\) be the same energy with
\(\mathbb C^c\) in place of \(\mathbb C\), and
\(\Pi^c=\partial\psi^c/\partial\bar U^c\).

\hypertarget{three-lemmas}{%
\section{3. Three lemmas}\label{three-lemmas}}

\textbf{Lemma 1 (rotated material law).}
\(\psi(\bar U)=\psi^c(\bar U^c)\) and \(\Pi=R_0\,\Pi^cR_0^T\).

\emph{Proof.} Under \(\bar U\mapsto R_0^T\bar UR_0\):

\begin{itemize}
\tightlist
\item
  \(J\) is unchanged, and \(\bar E\mapsto R_0^T\bar ER_0\).
\item
  \((R_0\star\mathbb C^c):\bigl(R_0AR_0^T\bigr)=R_0(\mathbb C^c:A)R_0^T\),
  so
  \(\tfrac12\bar E:\mathbb C:\bar E=\tfrac12\bar E^c:\mathbb C^c:\bar E^c\).
\item
  \(|a|\) is unchanged, since
  \(\operatorname{axl}(Q^TWQ)=Q^T\operatorname{axl}W\).
\end{itemize}

The relation for \(\Pi\) is the chain rule for the linear map
\(\bar U\mapsto R_0^T\bar UR_0\). \(\square\)

\textbf{Lemma 2 (no elastic torque).} \(\bar U^T\bar U=F_e^TF_e\) for
every \(R\). Hence \(\psi_{el}\) is independent of \(R\), and the
lattice rotation is set only by the constraint and curvature terms.

\textbf{Lemma 3 (the constraint selects the polar rotation).} Suppose
\(a=0\) and \(\bar U\succ0\). Then \(F_e=R\bar U\) is a polar
decomposition of \(F_e\), and by uniqueness
\[R=R_e=\operatorname{polar}(FF_p^{-1}),\qquad \bar U=U_e .\] The branch
\(\bar U\succ0\) is selected by continuity from
\(R(0)=I=\operatorname{polar}(F_e(0))\). (See the constraint derivation,
constraint\_finite.pdf.) \(\square\)

\hypertarget{theorem-the-stress-follows-the-lattice}{%
\section{4. Theorem: the stress follows the
lattice}\label{theorem-the-stress-follows-the-lattice}}

\textbf{Theorem 1.} For any \(F\) and \(F_p\),
\[\tau=R_t\,\mathcal T^c(\bar U^c)\,R_t^T,\qquad \mathcal T^c(U):=\Pi^c(U)\,U^T. \tag{1}\]
On the constraint (\(a=0\)), \(\bar U^c=U_e^c\), where
\[U_e^c=R_t^TF_eR_0=R_0^TU_eR_0 \tag{2}\] is the elastic stretch
measured in the lattice frame. Equation (1) then reads: \textbf{the
stress is the crystal's own response to its own elastic stretch, pushed
forward by the current lattice orientation.}

\emph{Proof.} 1.
\(\tau=PF^T=R\Pi F_p^{-T}F_p^T\bar U^TR^T=R\,\Pi\,\bar U^TR^T\). 2.
Lemma 1 gives \(\Pi=R_0\Pi^cR_0^T\) and \(\bar U=R_0\bar U^cR_0^T\), so
\[\tau=RR_0\,\Pi^c\bar U^{cT}R_0^TR^T=R_t\,\mathcal T^c(\bar U^c)\,R_t^T .\]
3. On the constraint, Lemma 3 gives \(\bar U=U_e\) and \(R=R_e\). Hence
\(\bar U^c=R_0^TU_eR_0=R_0^TR_e^TF_eR_0=R_t^TF_eR_0\). \(\square\)

\textbf{Corollary 1 (lattice covariance).} Consider two states with the
same lattice-frame stretch \(U_e^c\) and lattice orientations \(R_t\)
and \(R_t'\). Their stresses satisfy
\[\tau'=Q\,\tau\,Q^T,\qquad Q=R_t'R_t^T,\] independently of \(F\),
\(F_p\) and their histories.

\textbf{Corollary 2 (small elastic strain).} With
\(U_e^c=I+\varepsilon^c\) and the spatial strain
\(\varepsilon=R_t\varepsilon^cR_t^T\),
\[\tau\approx(R_t\star\mathbb C^c):\varepsilon .\] The spatial stiffness
is the crystal stiffness rotated by the current lattice orientation.

\textbf{Remark (the elastic stress follows \(F_e\) even without the
constraint).} For a Green-strain law \(\Pi_{el}=\bar U\,S\) with
\(S=S(\bar U^T\bar U)\), so
\[\tau_{el}=R\bar US\bar U^TR^T=F_e\,S(F_e^TF_e)\,F_e^T .\] This depends
only on \(F_e\), so the elastic stress always follows the lattice
carried by \(F_e\). What the constraint adds is that the Cosserat
rotation \(R\), which enters the curvature
\(K_i=\operatorname{axl}(R_t^T\partial_iR_t)\) and hence the
grain-boundary energy, is that same lattice.

\textbf{Remark (skew stress).} The multiplier and the penalty add
\(R\,(2\lambda+4\mu_c a)\,(\partial a/\partial\bar U)\,\bar U^TR^T\) to
\(\tau\); in 2D,
\(\partial a/\partial\bar U=\tfrac12(e_2\otimes e_1-e_1\otimes e_2)\).
This is the reaction to the curvature torque. It is expressed in the
same lattice frame, and it vanishes when there is no curvature torque.

\hypertarget{theorem-slip-follows-the-lattice}{%
\section{5. Theorem: slip follows the
lattice}\label{theorem-slip-follows-the-lattice}}

\textbf{Theorem 2.} On the constraint, with \(s_\alpha\) and
\(n_\alpha\) fixed in the intermediate configuration:

\begin{enumerate}
\def\labelenumi{\arabic{enumi}.}
\tightlist
\item
  \textbf{Driving force.} The resolved shear stress is the Kirchhoff
  stress resolved on the current lattice slip system:
  \[\tau_\alpha=M:(s_\alpha\otimes n_\alpha)=\tau:\bigl(\tilde s_\alpha\otimes\tilde n_\alpha\bigr),\qquad
  \tilde s_\alpha:=F_es_\alpha=R_tU_e^cs^c_\alpha,\quad \tilde n_\alpha:=F_e^{-T}n_\alpha=R_tU_e^{c\,-T}n^c_\alpha .\tag{3}\]
  Equivalently,
  \(\tau_\alpha=\Pi^c:(U_e^cs^c_\alpha\otimes n^c_\alpha)\), which is
  the crystal-frame expression.
\item
  \textbf{Kinematics.} Slip shears the material along the current slip
  direction on the current slip plane:
  \[F_eL_pF_e^{-1}=\sum_\alpha\dot\gamma_\alpha\,\tilde s_\alpha\otimes\tilde n_\alpha+\dot w^*\,F_e\widehat WF_e^{-1}. \tag{4}\]
\end{enumerate}

\emph{Proof.}

\begin{enumerate}
\def\labelenumi{(\arabic{enumi})}
\tightlist
\item
  We have \(M=\bar U^T\Pi\). Since \(F_e=R\bar U\) and
  \(\tau=R\Pi\bar U^TR^T\),
  \[F_e^T\tau F_e^{-T}=\bar U^TR^TR\Pi\bar U^TR^TR\bar U^{-T}=\bar U^T\Pi=M .\]
  Hence
  \[M:(s\otimes n)=s\cdot Mn=s\cdot F_e^T\tau F_e^{-T}n=(F_es)\cdot\tau\,(F_e^{-T}n)=\tau:(F_es\otimes F_e^{-T}n).\]
\end{enumerate}

The lattice-frame forms follow from \(F_e=R_tU_e^cR_0^T\) (Theorem 1,
Eq. 2) and \(s_\alpha=R_0s^c_\alpha\), \(n_\alpha=R_0n^c_\alpha\). The
crystal-frame expression follows from \(M=\bar U^T\Pi\) and Lemma 1.

\begin{enumerate}
\def\labelenumi{(\arabic{enumi})}
\setcounter{enumi}{1}
\tightlist
\item
  \(F_e(s\otimes n)F_e^{-1}=(F_es)\otimes(F_e^{-T}n)\). \(\square\)
\end{enumerate}

\textbf{Corollary 3.} The slip directions and normals must not be
rotated in time. Doing so in addition to the \(F_p\) update would
resolve the stress on a slip system rotated twice. The same holds for
\(\mathbb C\). (See elasticity\_invariance.pdf, Proposition 4, for the
equivalent ``rotating-lattice'' bookkeeping, in which \(F_p\) is updated
with slip only.)

\hypertarget{why-the-lattice-orientation-is-r_t-the-plastic-rotation}{%
\section{\texorpdfstring{6. Why the lattice orientation is \(R_t\): the
plastic
rotation}{6. Why the lattice orientation is R\_t: the plastic rotation}}\label{why-the-lattice-orientation-is-r_t-the-plastic-rotation}}

Theorem 2 of constraint\_finite.pdf gives, exactly,
\[R=\operatorname{polar}(FF_p^{-1})=R_F\,R_{\mathrm{mix}}\,R_p^T,\qquad R_{\mathrm{mix}}=\operatorname{polar}(U_FU_p^{-1}),\]
where \(F=R_FU_F\), \(F_p=R_pU_p\), and \(R_{\mathrm{mix}}=I\) for
coaxial stretches. In 2D this is
\(\theta=\theta_F-\theta_p+\theta_{\mathrm{mix}}\). For a pure plastic
rotation with \(F=I\), the lattice turns by \(R_p^T\) while the material
stays put (Figure 1).

\begin{figure}
\centering
\includegraphics[width=1\textwidth,height=\textheight]{lattice_frame_figs/polar_rotation_figure_ms.png}
\caption{Pure plastic rotation with \(F=I\). \(F_p=R(\theta_p)\) rotates
the material, and the lattice keeps its orientation in the intermediate
configuration. \(F_e=FF_p^{-1}=R(-\theta_p)\) returns the material and
carries the lattice. The constraint gives
\(R=\operatorname{polar}(FF_p^{-1})=R(-\theta_p)\). Red: material axes.
Blue: lattice.}
\end{figure}

The orientation \(R_t(X)\) is spatial: it says where the crystal axes at
material point \(X\) point now. The fields themselves are labelled by
the reference position \(X\).

\hypertarget{numerical-verification}{%
\section{7. Numerical verification}\label{numerical-verification}}

\textbf{Setup.}

\begin{itemize}
\tightlist
\item
  \textbf{Material.} Cu (\(C_{11}=168.4\), \(C_{12}=121.4\),
  \(C_{44}=75.4\) GPa), in the solver's lattice frame \(x=[110]\),
  \(y=[001]\), \(z=[\bar110]\). The setup reproduces the code's
  \(C_{xx}=220.3\), \(C_{yy}=168.4\), \(C_{xy}=121.4\) and
  \(C_{ss}=75.4\) GPa.
\item
  \textbf{Energy.} Plane-strain St Venant--Kirchhoff,
  \(\psi=\tfrac12E:\mathbb C:E+2\mu_ca^2\) with
  \(E=\tfrac12(\bar U^T\bar U-I)\), \(\mu_c=754\) GPa.
\item
  \textbf{Stress.} Production formulas: exact constraint
  \(R=\operatorname{polar}(F_e)\), multiplier from the torque balance,
  and \(P=R(\Pi+2\lambda\,\partial a/\partial\bar U)F_p^{-T}\).
  Derivatives by complex step (exact to round-off).
\item
  \textbf{Load set.} 35 stretch states (uniaxial, biaxial, compression,
  pure shear, 20 random large stretches, \(J\approx0.75\)--\(1.3\)),
  each at 12 lattice angles \(\alpha\) from \(-100^\circ\) to
  \(170^\circ\). The stress changes by up to 138\% of its own size as
  \(\alpha\) varies, so agreement is not trivial.
\item
  \textbf{Reference.} A crystal that starts at orientation \(\alpha\):
  \(\mathbb C=R(\alpha)\star\mathbb C^c\), slip systems rotated by
  \(\alpha\), \(F=V\), \(F_p=I\).
\end{itemize}

\textbf{Tests} (\(\mathbb C\) and the slip systems fixed at orientation
\(0\), lattice turned by \(\alpha\) during the history):

\begin{itemize}
\tightlist
\item
  \textbf{T1, rigid rotation:} \(F=VR(\alpha)\), \(F_p=I\).
\item
  \textbf{T2, grain-boundary spin:} \(F=V\), \(F_p=R(-\alpha)\).
\item
  \textbf{T3, slip plus spin:} \(F_p\) = three slips
  (\(\gamma=0.5,-0.2,0.3\) on systems at \(0^\circ\), \(60^\circ\),
  \(120^\circ\)) followed by \(R(-\alpha)\), with \(F=VR(\alpha)F_p\).
\item
  \textbf{T4, time-stepped spin:} 200 steps,
  \(F_p\leftarrow\exp(\Delta w^*\widehat W)F_p\), \(w^*\to-73^\circ\),
  while \(F\) goes from \(I\) to \(V\).
\item
  \textbf{W, wrong:} T2 with \(\mathbb C\) also rotated by \(\alpha\).
\end{itemize}

\begin{longtable}[]{@{}
  >{\raggedright\arraybackslash}p{(\columnwidth - 10\tabcolsep) * \real{0.1667}}
  >{\raggedright\arraybackslash}p{(\columnwidth - 10\tabcolsep) * \real{0.1667}}
  >{\raggedright\arraybackslash}p{(\columnwidth - 10\tabcolsep) * \real{0.1667}}
  >{\raggedright\arraybackslash}p{(\columnwidth - 10\tabcolsep) * \real{0.1667}}
  >{\raggedright\arraybackslash}p{(\columnwidth - 10\tabcolsep) * \real{0.1667}}
  >{\raggedright\arraybackslash}p{(\columnwidth - 10\tabcolsep) * \real{0.1667}}@{}}
\toprule\noalign{}
\begin{minipage}[b]{\linewidth}\raggedright
Quantity
\end{minipage} & \begin{minipage}[b]{\linewidth}\raggedright
T1
\end{minipage} & \begin{minipage}[b]{\linewidth}\raggedright
T2
\end{minipage} & \begin{minipage}[b]{\linewidth}\raggedright
T3
\end{minipage} & \begin{minipage}[b]{\linewidth}\raggedright
T4
\end{minipage} & \begin{minipage}[b]{\linewidth}\raggedright
W
\end{minipage} \\
\midrule\noalign{}
\endhead
\bottomrule\noalign{}
\endlastfoot
Kirchhoff stress \(\tau\) (max rel. error) & \(1.4\times10^{-14}\) &
\(1.9\times10^{-14}\) & \(1.6\times10^{-14}\) & \(2.4\times10^{-13}\) &
\textbf{1.70} \\
W compared with a crystal at \(2\alpha\) & & & & &
\(2.3\times10^{-14}\) \\
\end{longtable}

Slip systems, over T1--T3 and the full load set:

\begin{longtable}[]{@{}
  >{\raggedright\arraybackslash}p{(\columnwidth - 2\tabcolsep) * \real{0.5000}}
  >{\raggedright\arraybackslash}p{(\columnwidth - 2\tabcolsep) * \real{0.5000}}@{}}
\toprule\noalign{}
\begin{minipage}[b]{\linewidth}\raggedright
Quantity
\end{minipage} & \begin{minipage}[b]{\linewidth}\raggedright
max error
\end{minipage} \\
\midrule\noalign{}
\endhead
\bottomrule\noalign{}
\endlastfoot
Resolved shear stress \(\tau_\alpha\) (relative) &
\(3\times10^{-13}\) \\
Current slip direction \(F_es_\alpha\) & \(7\times10^{-16}\) \\
Current slip-plane normal \(F_e^{-T}n_\alpha\) & \(7\times10^{-16}\) \\
\(F_eL_pF_e^{-1}=\sum\dot\gamma_\alpha\,\tilde s_\alpha\otimes\tilde n_\alpha\)
& \(6\times10^{-17}\) \\
\end{longtable}

The penalty was also tested with an applied couple, so that \(a\neq0\)
and the penalty carries stress (\(|a|\) up to \(2.3\times10^{-2}\)).
T1--T3 still agree to \(\le6\times10^{-14}\); W does not.

\begin{figure}
\centering
\includegraphics[width=1\textwidth,height=\textheight]{lattice_frame_figs/c_fixed_all5.png}
\caption{Stress against lattice rotation for
\(V=\operatorname{diag}(1.10,0.95)\). Lines: reference crystal at
\(\alpha\). Markers: tests with \(\mathbb C\) fixed (T1--T4) and the
wrong variant (W, which falls on the curve for \(2\alpha\)). Last panel:
maximum error over the full load set.}
\end{figure}

\hypertarget{scope}{%
\section{8. Scope}\label{scope}}

\begin{itemize}
\tightlist
\item
  \textbf{Exact versus approximate.} The statements are exact when
  \(a=0\). In production the P0 multiplier makes the element mean of
  \(a\) zero exactly; within an element, the \(2\mu_ca^2\) penalty keeps
  \(|\theta-\theta_e|\) of order
  \((\text{curvature torque})/(\mu_c(\operatorname{tr}U_e)^2)\).
\item
  \textbf{Diffuse grain boundaries.} \(\mathbb C\), \(s_\alpha\) and
  \(n_\alpha\) use each grain's \(R_0\), taken from the mesh region. The
  curvature uses the smoothed \(\theta_0\). Inside the diffuse boundary
  the two initial orientations differ.
\item
  \textbf{The \(\pi\) branch.} A lattice rotated by \(\pi\) about an
  eigenvector of \(U_e\) gives an identical stress and identical
  \(\tau_\alpha\). ``The lattice'' is made unique by continuity from
  \(R(0)=I\), and in 2D by the symmetry of the square lattice.
\item
  \textbf{What the tests cover.} They are pointwise constitutive checks:
  no finite elements, no curvature, no time integration of the balance
  laws.
\end{itemize}

\textbf{Tests in the repository.}
\texttt{checks/test\_stress\_follows\_lattice.py} and
\texttt{checks/test\_slip\_follows\_lattice.py} (both laws, pass/fail,
figures); see \texttt{checks/LATTICE\_FRAME\_TESTS.md}.

\end{document}
